\section{Conventions and the induction hypotheses}\label{sec:preliminaries}

Until the final field-descent argument we work over $\C$. Varieties are
integral and projective unless another setting is stated. We use the
standard discrepancy conventions for klt, dlt, and lc pairs
\cite{Kollar2013}. All boundaries in the proof are rational. A
\emph{rational Cartier divisor} means a $\Q$-Cartier $\Q$-divisor.
Canonical divisors are chosen compatibly when forms are pulled back
or residues are taken. On a normal variety, sections of an integral
Weil divisor always mean sections of its rank-one reflexive divisorial
sheaf.

For rational divisors $D,E$ we write $D\preceq E$ if $E-D$ is rationally
linearly equivalent to an effective divisor. An index bound is a
\emph{common trivializing multiple}: a single integer kills all the
actual divisor classes under consideration and clears their
coefficients. A bound on the least possible integers implies a common
multiple by taking their least common multiple.

A \emph{contraction} is a projective surjective morphism
$f:X\to Z$ of normal varieties such that $f_*\OO_X=\OO_Z$. Its generic
fibre is geometrically integral in characteristic zero. A rational
fibration is understood through a relatively algebraically closed
subfield of the total function field; after resolution and Stein
factorization it has connected fibres. A smooth geometric generic
model is a smooth projective model of the geometric generic fibre.

For generalized pairs, the nef data consist of a nef rational Cartier
divisor on a higher birational model, together with its compatible
pullbacks. A moduli b-divisor is \emph{b-nef and b-Cartier} if it is the
Cartier closure of a nef divisor on one such model. The generalized
discrepancies are computed by pulling back the full adjoint, including
these nef data. We use the conventions of \cite{BirkarZhang2016}.

\begin{theorem}[Good minimal models; \cite{LA}]
\label{thm:good-models}
Let $(X,B)$ be a projective lc pair over $\C$ with effective rational
boundary and pseudo-effective $K_X+B$. Then $(X,B)$ has a good log
minimal model. In particular, its adjoint is nonvanishing, and on the
good model the adjoint is semiample.
\end{theorem}

This is the rational-boundary case of \cite[Theorem~11.1]{LA}.
We also use the standard klt MMP and its discrepancy comparisons
\cite{BCHM2010}. For klt pairs with a good minimal model, an MMP with
scaling can be chosen terminating at such a model. We will specify
relative or equivariant uses when they arise; they do not require a
separate unmentioned abundance assumption.

\begin{notation}[Induction statements]\label{not:induction}
Let $I_n$ denote Theorem~\ref{thm:main} in dimension $n$, and let $L_n$
denote Theorem~\ref{thm:lc-index} in dimension $n$, for every rational
DCC coefficient set. Let $K_n$ denote its zero-boundary klt case:
there is a common integer $a_n>0$ with $a_nK_X\sim0$ for all projective
klt $n$-folds satisfying $K_X\sim_\Q0$.
\end{notation}

The dimension-zero statements hold with degree one. In a fixed
induction step we assume $I_j,L_j$ for all $j<n$. By global ACC
\cite[Theorem~1.5]{HMX2014}, the coefficients of numerically trivial
lc pairs of fixed dimension with coefficients in a fixed DCC set
belong to a finite subset. Thus proving $L_n$ for every finite
rational coefficient set proves its stated DCC form. This observation
also applies to the different coefficients that occur in adjunction.

The dependency diagram below distinguishes $K_n$ from $L_n$. In
particular, the proof of the moduli denominator uses only $L_j$ with
$j<n$, not the dimension-$n$ index statement being established.
\begin{figure}[htbp]
\centering
\begin{tikzpicture}[
 box/.style={draw,rounded corners=2pt,align=center,font=\small,
             text width=39mm,minimum height=12mm,inner sep=4pt},
 arr/.style={-{Stealth[length=2mm]},semithick}]
\node[box] (lowL) at (0,0) {$L_j$, $j<n$\\normal lc indices};
\node[box] (lowI) at (6,0) {$I_j$, $j<n$\\effective Iitaka maps};
\node[box,text width=43mm] (den) at (0,-2) {Moduli denominators\\torsion over rationally\\connected bases};
\node[box] (scalar) at (6,-2) {Small volumes\\and scalar estimates};
\node[box,text width=57mm] (idx) at (6,-4) {$K_n$: zero-boundary klt index\\diagonal products and chains};
\node[box] (lc) at (6,-6) {$L_n$\\adjunction and Stein degree};
\node[box] (iit) at (0,-6) {$I_n$\\denominators and base birationality};
\draw[arr] (lowL)--(den);
\draw[arr] (lowI)--(scalar);
\draw[arr] (den)--(idx);
\draw[arr] (scalar)--(idx);
\draw[arr] (idx)--(lc);
\draw[arr] (lc)--(iit);
\draw[arr] (den)--(iit);
\end{tikzpicture}
\caption{One induction step. The left-hand route supplies the relative
index control used in reducing $K_n$; the right-hand route supplies its
section-order estimates. After $K_n$ is proved, adjunction and Stein
degree give $L_n$, and the denominator bound also feeds the final
proof of $I_n$, as shown by the direct left-hand arrow. Arrows record
the main dependencies, not morphisms of
varieties.}\label{fig:induction}
\end{figure}


\subsection{Geometric generic fibres and fields of definition}

All finite collections of varieties, maps, divisors, and rational
forms can be defined over a finitely generated field of characteristic
zero. Algebraically closed extensions preserve the singularity
conditions, dimensions of section spaces, Cartierness, and triviality
of the line bundles used here. For singularities one may spread a log
resolution and its discrepancy calculation; for sections and
trivializations one uses flat base change and faithful flatness.
Consequently the induction hypotheses apply to geometric generic
fibres, by comparison with an algebraically closed field embedded in
$\C$.

We occasionally need very general points to satisfy countably many
section-order conditions simultaneously. One may make these choices
over $\C$, or use points geometric generic over a countable field of
definition for the varieties and all the systems involved. Moving
members below are taken at geometric generic parameters. Dominant
extensions needed to define models or maps do not specialize a
previously chosen generic parameter.

We record one descent observation used repeatedly. If $F$ is a
geometrically integral normal projective variety over a field $K$
and a given reflexive pluricanonical sheaf becomes the trivial line
bundle over $\overline K$, then it is already a line bundle over $K$
and is trivial there. Local freeness descends faithfully flatly.
Base change identifies its space of sections with a one-dimensional
space after extension, and any nonzero section over $K$ becomes a
generator, hence is a generator before extension as well. The same
argument applies to a specified integral log pluricanonical divisor.
